% ============================================================
%  MANUAL DE USUARIO — ASESOR DE NEGOCIOS IA (Bot de Discord)
%  Sistema Experto · Cafetería de Especialidad GDL
% ============================================================
\documentclass[12pt, a4paper]{article}

% --- Idioma y codificación ---
\usepackage[spanish]{babel}
\decimalpoint
\usepackage[utf8]{inputenc}
\usepackage[T1]{fontenc}

% --- Tipografía ---
\usepackage{lmodern}
\usepackage{microtype}

% --- Márgenes ---
\usepackage[top=2.5cm, bottom=2.5cm, left=3cm, right=2.5cm]{geometry}

% --- Colores y gráficos ---
\usepackage[dvipsnames, table]{xcolor}
\usepackage{tikz}
\usetikzlibrary{arrows.meta, shapes.geometric, positioning, fit, backgrounds, calc, babel}
\usepackage{pgfplots}
\pgfplotsset{compat=1.18}

% --- Tablas ---
\usepackage{booktabs}
\usepackage{multirow}
\usepackage{array}
\usepackage{tabularx}
\usepackage{longtable}

% --- Matemáticas ---
\usepackage{amsmath, amssymb}

% --- Código ---
\usepackage{listings}
\lstset{
  basicstyle=\ttfamily\small,
  backgroundcolor=\color{gray!10},
  frame=single,
  rulecolor=\color{gray!40},
  breaklines=true,
  keywordstyle=\color{NavyBlue}\bfseries,
  commentstyle=\color{gray},
  stringstyle=\color{ForestGreen},
  language=Python,
}

% --- Cajas y alertas ---
\usepackage{tcolorbox}
\tcbuselibrary{skins, breakable}

\newtcolorbox{infobox}[1][]{
  colback=NavyBlue!8, colframe=NavyBlue!60,
  title=#1, fonttitle=\bfseries,
  arc=3pt, boxrule=0.8pt,
}
\newtcolorbox{warnbox}[1][]{
  colback=Orange!10, colframe=Orange!70,
  title=#1, fonttitle=\bfseries,
  arc=3pt, boxrule=0.8pt,
}
\newtcolorbox{rulebox}{
  colback=ForestGreen!8, colframe=ForestGreen!60,
  arc=3pt, boxrule=0.8pt,
}

% --- Encabezados y pie ---
\usepackage{fancyhdr}
\pagestyle{fancy}
\fancyhf{}
\fancyhead[L]{\small Asesor de Negocios IA — Bot de Discord}
\fancyhead[R]{\small \leftmark}
\fancyfoot[C]{\thepage}

% --- Hipervínculos ---
\usepackage[colorlinks=true, linkcolor=NavyBlue, urlcolor=NavyBlue]{hyperref}

% --- Enumeración ---
\usepackage{enumitem}

% --- Imagen (por si acaso) ---
\usepackage{graphicx}

% ============================================================
\title{
  \vspace{-1cm}
  {\Large\bfseries Asesor de Negocios IA}\\[0.4em]
  {\large Bot de Discord — Sistema Experto}\\[0.3em]
  {\normalsize Cafetería de Especialidad en la Zona Metropolitana de Guadalajara}\\[1em]
  \rule{\linewidth}{0.5pt}\\[0.5em]
  {\normalsize Manual de Usuario y Documentación Técnica}
}
\author{Proyecto SE}
\date{\today}

% ============================================================
\begin{document}

\maketitle
\tableofcontents
\newpage

% ============================================================
\section{Introducción}
% ============================================================

\textbf{Asesor de Negocios IA} es un simulador económico de una cafetería de especialidad
ubicada en la Zona Metropolitana de Guadalajara (ZMG), implementado como bot de Discord.
El sistema constituye un \textbf{sistema experto}: el dominio de conocimiento es la gestión
económica del negocio, y el simulador aplica ese conocimiento a través de reglas de inferencia
para guiar al usuario en sus decisiones.

\subsection{Objetivo pedagógico}

El proyecto demuestra los componentes fundamentales de los sistemas expertos:
\begin{itemize}
  \item \textbf{Base de conocimiento}: fórmulas económicas y reglas IF-THEN calibradas
        a la realidad de GDL (precios MXN, zonas, eventos locales).
  \item \textbf{Motor de inferencia}: evaluación de reglas sobre el contexto del negocio.
  \item \textbf{Interfaz de explicación}: el Agente Supervisor traduce las decisiones
        del sistema a lenguaje natural.
\end{itemize}

\subsection{Tecnologías empleadas}

\begin{center}
\begin{tabular}{ll}
  \toprule
  Componente & Tecnología \\
  \midrule
  Plataforma & Discord (slash commands) \\
  Lenguaje   & Python 3.11+ \\
  Framework  & discord.py 2.x \\
  Base de datos & SQLite 3 \\
  ORM        & Sin ORM (sqlite3 nativo) \\
  \bottomrule
\end{tabular}
\end{center}

% ============================================================
\section{Arquitectura del Sistema}
% ============================================================

El sistema sigue un esquema de \textbf{tres agentes} como requisito de la rúbrica.
Cada agente tiene una responsabilidad única y se comunica con los demás a través
de la base de datos compartida.

\subsection{Diagrama de agentes}

\begin{center}
\begin{tikzpicture}[
  agent/.style={
    draw=NavyBlue, fill=NavyBlue!12, rounded corners=6pt,
    minimum width=3.2cm, minimum height=1.4cm,
    text width=3cm, align=center, font=\small\bfseries
  },
  db/.style={
    draw=ForestGreen, fill=ForestGreen!12,
    cylinder, shape border rotate=90, aspect=0.3,
    minimum width=2.8cm, minimum height=1cm,
    font=\small\bfseries, align=center
  },
  user/.style={
    draw=Orange!80, fill=Orange!15, rounded corners=4pt,
    minimum width=2.4cm, minimum height=1cm,
    font=\small\bfseries, align=center
  },
  arr/.style={-{Stealth[length=7pt]}, thick},
  darr/.style={<->, thick, dashed, gray},
  node distance=1.6cm
]

% Nodos
\node[user] (U) {Usuario\\Discord};
\node[agent, right=2.2cm of U] (A1)
  {Agente 1\\Atención\\(AgenteAtencion)};
\node[agent, right=2cm of A1] (A2)
  {Agente 2\\Pedido/Inferencia\\(AgentePedido)};
\node[agent, below=1.8cm of A2] (A3)
  {Agente 3\\Supervisor\\(AgenteSupervisor)};
\node[db, below=1.8cm of A1] (DB) {Base de\\Datos SQLite};

% Flechas principales
\draw[arr] (U) -- node[above, font=\tiny]{comando} (A1);
\draw[arr] (A1) -- node[above, font=\tiny]{intención} (A2);
\draw[arr] (A2) -- node[right, font=\tiny]{persiste} (DB);
\draw[arr] (DB) -- node[left, font=\tiny]{lee} (A3);
\draw[arr] (A3) -| node[below, font=\tiny, near end]{explicación} (U);
\draw[darr] (A1) -- (DB);

% Etiquetas de BD
\node[font=\tiny, gray, below=0.05cm of DB, align=center, text width=4cm]
  {negocios $\cdot$ inventario $\cdot$ inferencias \\ finanzas $\cdot$ eventos $\cdot$ historial};

\end{tikzpicture}
\end{center}

\subsection{Descripción de los agentes}

\begin{infobox}[Agente 1 — Atención (\texttt{AgenteAtencion})]
Interpreta el comando o texto del usuario y lo convierte en una
\textbf{Intención estructurada} (\texttt{accion}, \texttt{parametros}).
Acciones posibles: \texttt{avanzar\_ciclo}, \texttt{comprar},
\texttt{ajustar\_precio}, \texttt{desconocida}.
\end{infobox}

\begin{infobox}[Agente 2 — Pedido/Inferencia (\texttt{AgentePedido})]
Núcleo del sistema. Ejecuta el siguiente pipeline por ciclo:
\begin{enumerate}[leftmargin=*]
  \item Construye el contexto económico del negocio (zona, inventario, marketing).
  \item Inyecta el evento macroeconómico global del ciclo.
  \item Calcula la demanda con elasticidad-precio.
  \item Corre el motor de reglas IF-THEN.
  \item Resuelve el ciclo económico (ventas, costos, flujo).
  \item Persiste todo en la BD.
  \item Consume inventario (con merma 7.5\%) y decae marketing.
\end{enumerate}
\end{infobox}

\begin{infobox}[Agente 3 — Supervisor (\texttt{AgenteSupervisor})]
Lee \texttt{registro\_inferencias} y traduce las decisiones del motor a
\textbf{lenguaje natural}. Provee la explicabilidad requerida por los
sistemas expertos. El usuario lo activa con el botón \textbf{Ver explicación}.
\end{infobox}

\subsection{Flujo completo de un ciclo}

\begin{center}
\begin{tikzpicture}[
  box/.style={draw=NavyBlue!70, fill=NavyBlue!10, rounded corners=4pt,
              minimum width=3.4cm, minimum height=0.8cm,
              font=\small, align=center},
  dec/.style={draw=Orange!80, fill=Orange!10, diamond, aspect=2.2,
              font=\scriptsize, align=center},
  arr/.style={-{Stealth[length=6pt]}, thick},
  node distance=0.85cm
]

\node[box] (start) {Usuario: /ciclo\_avanzar};
\node[box, below=of start] (ctx)   {Construir contexto\\(zona, inventario, precios)};
\node[box, below=of ctx]   (ev)    {Generar evento global\\(CATALOGO aleatorio)};
\node[box, below=of ev]    (dem)   {Calcular demanda\\(elasticidad-precio)};
\node[box, below=of dem]   (rules) {Motor de reglas\\IF-THEN};
\node[box, below=of rules] (sim)   {Resolver ciclo\\(ventas, costos, flujo)};
\node[box, below=of sim]   (pers)  {Persistir en BD\\(finanzas, evento, inferencias)};
\node[box, below=of pers]  (inv)   {Consumir inventario\\(merma 7.5\%)};
\node[box, below=of inv]   (mkt)   {Decaer marketing $\times 0.5$};
\node[box, below=of mkt]   (out)   {Embed resultado\\+ Panel de botones};

\foreach \a/\b in {start/ctx, ctx/ev, ev/dem, dem/rules,
                    rules/sim, sim/pers, pers/inv, inv/mkt, mkt/out}
  \draw[arr] (\a) -- (\b);

\end{tikzpicture}
\end{center}

% ============================================================
\section{Motor de Inferencias}
% ============================================================

\subsection{Arquitectura del motor}

El motor evalúa una lista ordenada de funciones \texttt{Regla :: dict -> Inferencia | None}.
Cada regla es \textbf{pura}: no modifica la BD, solo razona sobre el contexto.
El Agente 2 ejecuta las reglas y persiste los resultados en \texttt{registro\_inferencias}.

\begin{lstlisting}[caption={Estructura de una regla}]
def r_precio_alto(ctx: dict) -> Inferencia | None:
    if ctx["precio_usuario"] > ctx["precio_mercado"]:
        sobreprecio = (ctx["precio_usuario"] / ctx["precio_mercado"] - 1) * 100
        return Inferencia(
            regla="IF precio > mercado THEN demanda cae exponencialmente",
            entrada=f"precio={ctx['precio_usuario']:.0f}, mercado={ctx['precio_mercado']:.0f}",
            resultado=f"Sobreprecio de {sobreprecio:.0f}% -> menos clientes",
            explicacion="Cobras por encima del promedio...",
        )
    return None
\end{lstlisting}

\subsection{Catálogo de reglas}

\begin{center}
\begin{tabularx}{\linewidth}{>{\ttfamily\small}l X X}
  \toprule
  \textbf{Regla} & \textbf{Condición (IF)} & \textbf{Acción (THEN)} \\
  \midrule
  r\_precio\_alto &
    \texttt{precio\_usuario > precio\_mercado (\$55)} &
    Dispara alerta de sobreprecio; la elasticidad castiga la demanda exponencialmente. \\
  \addlinespace
  r\_stock\_insuficiente &
    \texttt{inventario\_tazas < demanda\_estimada} &
    Alerta de desabasto: se pierden ventas. Sugiere reabastecer con \texttt{/proveedor}. \\
  \addlinespace
  r\_grano\_especialidad &
    \texttt{tipo\_grano == "especialidad"} &
    +15\% reputación/demanda. La calidad del grano compensa precios altos. \\
  \addlinespace
  r\_riesgo\_bancarrota &
    \texttt{caja < costos\_fijos} &
    \textbf{Alerta crítica}: la caja no cubre los costos fijos del ciclo. \\
  \bottomrule
\end{tabularx}
\end{center}

\subsection{Estructura de datos \texttt{Inferencia}}

\begin{lstlisting}
@dataclass
class Inferencia:
    regla:       str   # nombre legible de la regla disparada
    entrada:     str   # valores que la activaron
    resultado:   str   # que se decidio
    explicacion: str   # lenguaje natural para el Supervisor
    efectos:     dict  # cambios a aplicar (ej. reputacion_mult: 1.15)
\end{lstlisting}

% ============================================================
\section{Modelo Económico}
% ============================================================

\subsection{Fórmulas principales}

\subsubsection{Demanda con elasticidad-precio}

\begin{equation}
  D = D_{\text{base}} \cdot \left(\frac{P_{\text{mercado}}}{P_{\text{usuario}}}\right)^{1.5}
\end{equation}

donde $P_{\text{mercado}} = \$55$ MXN/taza y el exponente $1.5$ modela una
elasticidad moderada-alta (especialidad en GDL es sensible al precio).

\subsubsection{Unidades vendidas}

\begin{equation}
  U = \min\bigl(D,\ C_{\text{capacidad}},\ I_{\text{inventario}}\bigr)
\end{equation}

con $C_{\text{capacidad}} = 700$ tazas/ciclo.

\subsubsection{Margen de contribución}

\begin{equation}
  MC = P_{\text{usuario}} - CV_{\text{unit}}
\end{equation}

\subsubsection{Punto de equilibrio}

\begin{equation}
  PE = \frac{CF}{MC} \quad \text{(tazas por ciclo)}
\end{equation}

\subsubsection{Flujo neto}

\begin{equation}
  FN = U \cdot P_{\text{usuario}} - \bigl(U \cdot CV_{\text{unit}} + CF\bigr)
\end{equation}

\subsubsection{Costo variable unitario con merma}

\begin{equation}
  CV_{\text{unit}} = \bigl(C_{\text{grano}} + C_{\text{leche}} + C_{\text{vaso}}\bigr) \times (1 + m)
\end{equation}

donde $m = 7.5\%$ es la merma estimada por ciclo.

\subsection{Parámetros calibrados a GDL (MXN)}

\begin{center}
\begin{tabular}{llr}
  \toprule
  \textbf{Parámetro} & \textbf{Descripción} & \textbf{Valor} \\
  \midrule
  \multicolumn{3}{l}{\textit{Costos variables por taza}} \\
  $C_{\text{grano comercial}}$   & Grano estándar         & \$6.00 MXN \\
  $C_{\text{grano especialidad}}$& Chiapas/Veracruz/Nayarit & \$14.00 MXN \\
  $C_{\text{leche}}$             & Leche entera           & \$4.00 MXN \\
  $C_{\text{vegetal}}$           & Avena/almendra         & \$9.00 MXN \\
  $C_{\text{vaso}}$              & Empaque desechable     & \$3.00 MXN \\
  $m$                            & Merma                  & 7.5\% \\
  \addlinespace
  \multicolumn{3}{l}{\textit{Costos fijos por ciclo (mensual) — base}} \\
  Renta base                     & $\times$ mult. zona    & \$18{,}000 MXN \\
  Barista Jr.                    & Nómina                 & \$8{,}500 MXN \\
  Servicios (CFE/agua/gas)       &                        & \$4{,}500 MXN \\
  Mantenimiento espresso         &                        & \$1{,}500 MXN \\
  \textbf{Total base}            &                        & $\mathbf{\$32{,}500}$ \textbf{MXN} \\
  \addlinespace
  \multicolumn{3}{l}{\textit{Otros parámetros}} \\
  $P_{\text{mercado}}$           & Precio referencia GDL  & \$55 MXN/taza \\
  Precio inicial                 & Al crear negocio       & \$55 MXN/taza \\
  Caja inicial                   &                        & \$100{,}000 MXN \\
  Capacidad                      & Tazas/ciclo            & 700 tazas \\
  Demanda base                   & $\times$ mult. zona    & 500 tazas \\
  \bottomrule
\end{tabular}
\end{center}

\subsubsection{Ejemplo: costos variables calculados}

\begin{center}
\begin{tabular}{lcc}
  \toprule
  Combinación & $CV_{\text{unit}}$ (sin merma) & $CV_{\text{unit}}$ (con 7.5\% merma) \\
  \midrule
  Grano comercial + leche   & \$13.00 & \$13.98 \\
  Grano comercial + vegetal & \$18.00 & \$19.35 \\
  Especialidad + leche      & \$21.00 & \$22.58 \\
  Especialidad + vegetal    & \$26.00 & \$27.95 \\
  \bottomrule
\end{tabular}
\end{center}

\subsection{Sistema de zonas — ZMG}

La zona elegida al crear el negocio afecta la renta y la demanda base
mediante multiplicadores independientes:

\begin{equation}
  CF_{\text{zona}} = \$18{,}000 \times m_{\text{renta}} + \$14{,}500
\end{equation}
\begin{equation}
  D_{\text{base,zona}} = 500 \times m_{\text{demanda}}
\end{equation}

\begin{center}
\begin{tabular}{llccr}
  \toprule
  \textbf{Zona} & \textbf{Perfil} & $m_{\text{renta}}$ & $m_{\text{demanda}}$ & \textbf{CF/ciclo} \\
  \midrule
  Centro Histórico      & Alto flujo, turismo          & 1.0 & 1.1 & \$32{,}500 \\
  Col. Americana / Chap.& Hub de especialidad          & 1.4 & 1.3 & \$39{,}700 \\
  Providencia           & Comercial, poder adquisitivo & 1.5 & 1.2 & \$41{,}500 \\
  Andares / Pta. Hierro & Premium, plazas              & 1.8 & 1.1 & \$46{,}900 \\
  Tlaquepaque / Tonalá  & Turístico, artesanal         & 1.1 & 1.0 & \$34{,}300 \\
  Periferia             & Renta baja, menos tráfico    & 0.6 & 0.8 & \$25{,}300 \\
  \bottomrule
\end{tabular}
\end{center}

\subsubsection{Análisis de rentabilidad por zona}

El siguiente gráfico muestra el punto de equilibrio (en tazas/ciclo) para
cada zona, asumiendo precio \$55, grano comercial + leche:

\begin{center}
\begin{tikzpicture}
\begin{axis}[
  ybar, bar width=14pt,
  width=12cm, height=7cm,
  symbolic x coords={Periferia, Centro, Tlaquepaque,
                     Americana, Providencia, Andares},
  xtick=data,
  x tick label style={font=\small, rotate=20, anchor=east},
  ylabel={Punto de equilibrio (tazas/ciclo)},
  ymin=0, ymax=1400,
  ymajorgrids=true, grid style=dashed,
  nodes near coords,
  nodes near coords style={font=\tiny},
  bar shift=0pt,
  title={\small Punto de equilibrio por zona (precio \$55, grano comercial + leche)},
  title style={font=\small},
  cycle list={
    {fill=ForestGreen!70, draw=ForestGreen!90},
    {fill=NavyBlue!50,    draw=NavyBlue!70},
    {fill=NavyBlue!65,    draw=NavyBlue!85},
    {fill=Orange!65,      draw=Orange!85},
    {fill=Orange!80,      draw=Orange!100},
    {fill=red!65,         draw=red!85},
  },
]
\addplot coordinates {
  (Periferia,    613)
  (Centro,       792)
  (Tlaquepaque,  835)
  (Americana,    966)
  (Providencia, 1010)
  (Andares,     1141)
};
\end{axis}
\end{tikzpicture}
\end{center}

% ============================================================
\section{Base de Datos}
% ============================================================

\subsection{Diagrama entidad-relación}

\begin{center}
\begin{tikzpicture}[
  tbl/.style={draw=NavyBlue!70, fill=NavyBlue!8, rounded corners=4pt,
              font=\small, inner sep=6pt},
  hdr/.style={fill=NavyBlue!50, text=white, font=\small\bfseries, rounded corners=2pt,
              inner sep=4pt},
  rel/.style={-{Stealth[length=6pt]}, thick, NavyBlue!70},
  node distance=1.2cm and 2.2cm
]

% negocios
\node[tbl] (N) {
  \begin{tabular}{l}
    \rowcolor{NavyBlue!50}
    \color{white}\textbf{negocios} \\
    \underline{negocio\_id} \\
    owner\_id \\
    owner\_name \\
    nombre \\
    caja \\
    precio\_taza \\
    ciclo\_actual \\
    zona \\
    marketing\_bonus \\
  \end{tabular}
};

% inventario
\node[tbl, right=2.5cm of N] (INV) {
  \begin{tabular}{l}
    \rowcolor{ForestGreen!50}
    \color{white}\textbf{inventario} \\
    \underline{id} \\
    negocio\_id (FK) \\
    tipo\_insumo \\
    tazas\_equiv \\
  \end{tabular}
};

% finanzas_registro
\node[tbl, below=1.5cm of N] (FIN) {
  \begin{tabular}{l}
    \rowcolor{Orange!60}
    \color{white}\textbf{finanzas\_registro} \\
    \underline{id} \\
    negocio\_id (FK) \\
    ciclo \\
    ingresos \\
    costos\_fijos \\
    costos\_variables \\
    flujo\_neto \\
    caja\_final \\
  \end{tabular}
};

% historial_eventos
\node[tbl, right=2.5cm of FIN] (EV) {
  \begin{tabular}{l}
    \rowcolor{Orange!60}
    \color{white}\textbf{historial\_eventos} \\
    \underline{id} \\
    negocio\_id (FK) \\
    ciclo \\
    tipo \\
    descripcion \\
    impacto\_json \\
  \end{tabular}
};

% registro_inferencias
\node[tbl, below=1.5cm of FIN] (INF) {
  \begin{tabular}{l}
    \rowcolor{Plum!50}
    \color{white}\textbf{registro\_inferencias} \\
    \underline{id} \\
    negocio\_id (FK) \\
    ciclo \\
    agente \\
    regla \\
    entrada \\
    resultado \\
    explicacion \\
  \end{tabular}
};

% Relaciones
\draw[rel] (INV.west)  -- node[above, font=\tiny]{1:N} (N.east);
\draw[rel] (FIN.north) -- node[right, font=\tiny]{1:N} (N.south);
\draw[rel] (EV.west)   -- (FIN.east);
\draw[rel] (INF.north) -- node[right, font=\tiny]{1:N} (FIN.south);

\end{tikzpicture}
\end{center}

\subsection{Descripción de tablas}

\begin{longtable}{lll}
  \toprule
  \textbf{Tabla} & \textbf{Campo} & \textbf{Descripción} \\
  \midrule
  \endhead
  \texttt{negocios} & negocio\_id & PK autoincrement \\
  & owner\_id & Discord user ID \\
  & caja & Efectivo disponible (MXN) \\
  & precio\_taza & Precio de venta actual \\
  & zona & Zona ZMG elegida \\
  & marketing\_bonus & Multiplicador activo (decae 50\%/ciclo) \\
  \addlinespace
  \texttt{inventario} & tipo\_insumo & grano\_comercial, grano\_especialidad, leche, vegetal \\
  & tazas\_equiv & Tazas que puede producir el insumo disponible \\
  \addlinespace
  \texttt{finanzas\_registro} & ciclo & Número de ciclo (mes) \\
  & flujo\_neto & Ingresos $-$ (CF + CV) \\
  & caja\_final & Caja después del ciclo \\
  \addlinespace
  \texttt{historial\_eventos} & tipo & Categoría del evento (clima, tendencia, etc.) \\
  & impacto\_json & Multiplicadores aplicados en JSON \\
  \addlinespace
  \texttt{registro\_inferencias} & agente & Quién disparó la regla (pedido/atencion) \\
  & regla & Nombre legible de la regla IF-THEN \\
  & explicacion & Texto para el Agente Supervisor \\
  \bottomrule
\end{longtable}

% ============================================================
\section{Eventos Macroeconómicos}
% ============================================================

Cada ciclo el sistema elige \textbf{un evento} del catálogo mediante selección
ponderada. El evento ``Mes estable'' tiene peso $3\times$ para que el mercado
no sea caótico (peso $1$ para cada evento perturbador).

\begin{center}
\begin{tabular}{llll}
  \toprule
  \textbf{Tipo} & \textbf{Descripción} & \textbf{Variable afectada} & \textbf{Impacto} \\
  \midrule
  \texttt{clima}      & Helada en Brasil; grano sube       & \texttt{costo\_variable\_unit} & $\times 1.25$ \\
  \texttt{tendencia}  & Viral en TikTok; demanda sube      & \texttt{demanda\_base}         & $\times 1.40$ \\
  \texttt{inflacion}  & Inflación local; costos fijos suben& \texttt{costos\_fijos}         & $\times 1.05$ \\
  \texttt{competencia}& Rival cerca; demanda baja          & \texttt{demanda\_base}         & $\times 0.85$ \\
  \texttt{cosecha}    & Cosecha Colombia; grano baja       & \texttt{costo\_variable\_unit} & $\times 0.90$ \\
  \texttt{estable}    & Mes tranquilo                      & ---                            & sin cambio    \\
  \bottomrule
\end{tabular}
\end{center}

\begin{center}
\begin{tikzpicture}
\begin{axis}[
  ybar, bar width=16pt,
  width=10cm, height=5.5cm,
  symbolic x coords={clima, tendencia, inflacion, competencia, cosecha, estable},
  xtick=data,
  x tick label style={font=\small, rotate=25, anchor=east},
  ylabel={Peso de selección},
  ymin=0, ymax=4,
  ymajorgrids=true, grid style=dashed,
  nodes near coords,
  title={\small Distribución de probabilidad de eventos},
  title style={font=\small},
  cycle list={
    {fill=red!60, draw=red!80},
    {fill=green!60, draw=green!80},
    {fill=Orange!70, draw=Orange!90},
    {fill=red!40, draw=red!60},
    {fill=green!40, draw=green!60},
    {fill=gray!50, draw=gray!70},
  },
]
\addplot coordinates {
  (clima,1)(tendencia,1)(inflacion,1)(competencia,1)(cosecha,1)(estable,3)
};
\end{axis}
\end{tikzpicture}
\end{center}

% ============================================================
\section{Sistema de Inventario}
% ============================================================

El inventario almacena insumos como \textbf{tazas-equivalente}: la cantidad de
tazas que se pueden producir con el insumo disponible. Se descuenta
automáticamente al vender, con una merma del 7.5\%.

\subsection{Insumos disponibles}

\begin{center}
\begin{tabular}{lrrl}
  \toprule
  \textbf{Insumo} & \textbf{Costo/lote} & \textbf{Tazas/lote} & \textbf{Efecto} \\
  \midrule
  Grano comercial       & \$300 MXN  & 50 & Costo variable base \\
  Grano de especialidad & \$700 MXN  & 50 & +15\% demanda, mayor CV \\
  Leche entera          & \$120 MXN  & 30 & Insumo base \\
  Alternativas vegetales& \$270 MXN  & 30 & Mayor CV, segmento premium \\
  \bottomrule
\end{tabular}
\end{center}

\subsection{Flujo del inventario}

\begin{center}
\begin{tikzpicture}[
  box/.style={draw=NavyBlue!60, fill=NavyBlue!8, rounded corners=3pt,
              minimum width=3.2cm, minimum height=0.7cm, font=\small, align=center},
  arr/.style={-{Stealth[length=6pt]}, thick},
  node distance=0.8cm
]
\node[box] (a) {Crear negocio};
\node[box, right=1.5cm of a] (b) {Inventario inicial\\(100tz grano + 60tz leche)};
\node[box, right=1.5cm of b] (c) {Compra /proveedor\\(suma tazas)};
\node[box, below=1.2cm of c] (d) {Ciclo / Venta rápida\\(consume $U \times 1.075$)};
\node[box, below=1.2cm of b] (e) {Inventario actualizado};

\draw[arr] (a) -- (b);
\draw[arr] (b) -- (c);
\draw[arr] (c) -- (d);
\draw[arr] (d) -- (e);
\draw[arr] (e) -- (b) node[midway, left, font=\tiny]{reabastecer};
\end{tikzpicture}
\end{center}

% ============================================================
\section{Sistema de Marketing}
% ============================================================

El marketing se modela como un \textbf{bonus multiplicativo} sobre la demanda base
que decae un 50\% cada ciclo, incentivando la reinversión continua.

\begin{equation}
  D_{\text{base,final}} = D_{\text{base}} \times (1 + b_{\text{mkt}})
\end{equation}

\begin{equation}
  b_{\text{mkt}}^{(t+1)} = b_{\text{mkt}}^{(t)} \times 0.5
\end{equation}

\begin{center}
\begin{tabular}{lrrl}
  \toprule
  \textbf{Campaña} & \textbf{Costo} & \textbf{Bonus demanda} & \textbf{Contexto} \\
  \midrule
  Redes sociales            & \$3{,}000 MXN & +10\% & Bajo costo, efecto moderado \\
  Colaboración con influencer& \$8{,}000 MXN & +20\% & Mayor alcance \\
  Descuento de apertura     & \$1{,}500 MXN & +25\% & Más demanda, precio efectivo menor \\
  Evento de cata en tienda  & \$5{,}000 MXN & +15\% & Refuerza reputación de especialidad \\
  \bottomrule
\end{tabular}
\end{center}

\begin{center}
\begin{tikzpicture}
\begin{axis}[
  width=10cm, height=5cm,
  xlabel={Ciclos transcurridos},
  ylabel={Bonus activo (\%)},
  ymin=0, ymax=30,
  xmin=0, xmax=6,
  xtick={0,1,2,3,4,5,6},
  ymajorgrids=true, grid style=dashed,
  legend pos=north east,
  title={\small Decaimiento del bonus de marketing por campaña},
  title style={font=\small},
]
\addplot[NavyBlue, thick, mark=*]
  coordinates {(0,20)(1,10)(2,5)(3,2.5)(4,1.25)(5,0.6)(6,0.3)};
\addlegendentry{Influencer (+20\%)}

\addplot[ForestGreen, thick, mark=square*]
  coordinates {(0,10)(1,5)(2,2.5)(3,1.25)(4,0.6)(5,0.3)(6,0.15)};
\addlegendentry{Redes sociales (+10\%)}

\addplot[Orange, thick, mark=triangle*]
  coordinates {(0,25)(1,12.5)(2,6.25)(3,3.12)(4,1.56)(5,0.78)(6,0.39)};
\addlegendentry{Descuento (+25\%)}
\end{axis}
\end{tikzpicture}
\end{center}

% ============================================================
\section{Referencia de Comandos}
% ============================================================

\begin{center}
\begin{tabularx}{\linewidth}{llX}
  \toprule
  \textbf{Comando} & \textbf{Parámetros} & \textbf{Descripción} \\
  \midrule
  \texttt{/crear\_negocio} & \texttt{nombre}, \texttt{zona} &
    Crea la cafetería. Elige zona de 6 opciones de la ZMG. Caja inicial: \$100{,}000 MXN. \\
  \addlinespace
  \texttt{/ciclo\_avanzar} & --- &
    Simula el siguiente mes: evento, inferencias, cálculo económico, actualiza BD. \\
  \addlinespace
  \texttt{/dashboard} & --- &
    Panel completo: finanzas, inventario, último ciclo e indicadores. Selector si hay varios negocios. \\
  \addlinespace
  \texttt{/proveedor} & --- &
    Abre Select Menu para comprar insumos (descuenta caja, suma inventario). \\
  \addlinespace
  \texttt{/marketing} & --- &
    Abre Select Menu de 4 campañas de marketing. El bonus decae 50\%/ciclo. \\
  \addlinespace
  \texttt{/ajustar\_precio} & \texttt{precio} &
    Cambia el precio de venta. Valida que sea $\geq$ costo variable (\$13 MXN). \\
  \addlinespace
  \texttt{/borrar\_negocio} & --- &
    Elimina el negocio y todos sus datos con confirmación doble. \\
  \bottomrule
\end{tabularx}
\end{center}

\subsection{Componentes interactivos (botones y selectores)}

Tras ejecutar \texttt{/ciclo\_avanzar} aparece el \textbf{Panel de Ciclo}
con cuatro botones:

\begin{center}
\begin{tabular}{lll}
  \toprule
  \textbf{Botón} & \textbf{Color} & \textbf{Función} \\
  \midrule
  Avanzar ciclo   & Azul (primario) & Corre el siguiente mes completo. \\
  Ver explicación & Gris (secundario) & Muestra el reporte del Agente Supervisor. \\
  Vender          & Verde & Venta rápida sin avanzar ciclo ni aplicar costos fijos. \\
  Pagar Nómina    & Rojo & Descuenta \$8{,}500 MXN de caja manualmente. \\
  \bottomrule
\end{tabular}
\end{center}

% ============================================================
\section{Guía de Inicio Rápido}
% ============================================================

\begin{enumerate}[leftmargin=*, label=\textbf{\arabic*.}]
  \item \textbf{Crea tu cafetería}\\
        \texttt{/crear\_negocio nombre:"La Brújula" zona:"Col. Americana / Chapultepec"}\\
        Recibirás \$100{,}000 MXN de caja inicial y un inventario base de 100 tazas
        de grano comercial y 60 de leche.

  \item \textbf{Consulta tu estado}\\
        \texttt{/dashboard} — revisa tu caja, inventario actual e indicadores
        (margen de contribución y punto de equilibrio).

  \item \textbf{Avanza tu primer ciclo}\\
        \texttt{/ciclo\_avanzar} — el sistema generará un evento macroeconómico,
        calculará ventas y actualizará tu caja. Aparecerá el panel de botones.

  \item \textbf{Analiza las inferencias}\\
        Pulsa \textbf{Ver explicación} para que el Agente Supervisor te explique
        las reglas que se dispararon y por qué.

  \item \textbf{Reabastece y estrategiza}
  \begin{itemize}
    \item \texttt{/proveedor} → compra grano de especialidad para mejorar reputación.
    \item \texttt{/marketing} → lanza una campaña para subir la demanda.
    \item \texttt{/ajustar\_precio} → experimenta con el precio y observa la elasticidad.
  \end{itemize}
\end{enumerate}

\begin{warnbox}[Consejo estratégico]
La zona \textbf{Col. Americana / Chapultepec} ofrece el mejor balance
demanda/renta para un negocio rentable desde el ciclo 1. Las zonas
premium (Andares) requieren volumen alto para superar el punto de equilibrio.
\end{warnbox}

% ============================================================
\section{Ejemplo de Sesión}
% ============================================================

\begin{lstlisting}[language={}, caption={Ejemplo de salida de /ciclo\_avanzar}]
Ciclo 3 - La Brujula
Evento global: Tendencia viral en TikTok: la demanda sube 40%.

Tazas vendidas    490
Ingresos          $26,950
Flujo neto        +$4,543
Caja final        $71,234
Reglas disparadas  1
\end{lstlisting}

\begin{lstlisting}[language={}, caption={Reporte del Supervisor (Ver explicacion)}]
Decisiones tomadas por el sistema:

- IF grano = especialidad THEN reputacion +15%
    -> El grano de especialidad mejora la percepcion de calidad
       y la reputacion de la marca.
\end{lstlisting}

% ============================================================
\section{Consideraciones Técnicas}
% ============================================================

\subsection{Persistencia y transacciones}

Todas las operaciones sobre la BD se realizan con \texttt{conn.commit()}
explícito. Las operaciones críticas (compra de insumos, venta) verifican
fondos antes de ejecutar para garantizar consistencia.

\subsection{Pureza del motor de reglas}

Las funciones en \texttt{rules.py} son \textbf{puras}: reciben un contexto
como diccionario y devuelven una \texttt{Inferencia} o \texttt{None}, sin
efectos secundarios. Esto facilita pruebas unitarias y la explicabilidad.

\subsection{Escalabilidad del catálogo de reglas}

Añadir una nueva regla consiste en:
\begin{enumerate}
  \item Definir \texttt{def r\_nueva(ctx: dict) -> Inferencia | None} en \texttt{rules.py}.
  \item Añadirla a la lista \texttt{REGLAS}.
\end{enumerate}
No se requiere modificar ningún otro módulo.

\subsection{Tests sugeridos (pytest)}

\begin{lstlisting}[caption={Pruebas recomendadas para la presentación}]
def test_punto_equilibrio():
    cf, cv, p = 32500, 13.98, 55.0
    pe = cf / (p - cv)
    assert abs(pe - 793) < 5

def test_elasticidad_precio():
    d = demanda_estimada(500, 55, 110)  # precio doble
    assert d < 500 * 0.5  # demanda cae mas del 50%

def test_regla_bancarrota():
    ctx = {"caja": 1000, "costos_fijos": 32500, ...}
    inf = r_riesgo_bancarrota(ctx)
    assert inf is not None
\end{lstlisting}

\vfill
\hrule
\vspace{0.4em}
\begin{center}
  \small Prototipo Figma: \url{https://gown-kinder-24829009.figma.site/}
  \quad|\quad
  Proyecto SE — Sistemas Expertos
\end{center}

\end{document}
